% !TeX root = ../paper.tex
\subsection{Data Collection}
\label{sec:data-collection}

\subsubsection{Data Sources}
% Describe all data sources used in your project
% Be specific about where each type of data comes from

We collect our data from various authoritative and reputable sources to ensure the factual correctness of our model,
among which are textbooks, journals, and fact-sheets from worldwide institution.
Specifically, we collect textbooks and fact-sheets to build the knowledge base for \ac{RAG}, while for instruction fine-tuning,
we utilize the multi-modal \ac{CoT} reasoning QA-set from LangAGI-Lab, which contains sampled questions from \ac{USMLE}, 
a well-known medical licensing examination in the United States. Finally, for benchmarking, 
we synthesize a QA-set from 93 clinical case reports provided by the lecturer and 80 open-access case reports from the \ac{AJTMH}. 
The data sources in our collection are summarized in Table~\ref{tab:data-sources-type} and Table~\ref{tab:data-sources}.

% Add table summarizing data sources
\begin{table}[htbp]
\centering
\begin{tabular}{lll}
\hline
\textbf{Name} & \textbf{Source Type} & \textbf{Volume} \\
\hline
Manson's Tropical Infectious Diseases (23rd Edition) & Textbook & 80 chapters - 1655 pages \\
Manson's Tropical Diseases (24th Edition) & Textbook & 84 chapters - 1634 pages \\
Tropical and Parasitic Infections in the Intensive Care Unit & Textbook & 12 chapters - 247 pages \\
Current Topics in Tropical Medicine & Textbook & 29 chapters - 578 pages \\
WHO & Fact-sheet & 50 sheets \\
LangAGI-Lab & QA-set & 160 pairs \\
Moodle & Case Report & 93 articles \\
\acs{AJTMH} & Case Report & 80 articles \\
\hline
\end{tabular}
\caption{Overview of Data Sources}
\label{tab:data-sources}
\end{table}

\begin{table}[htbp]
\centering
\begin{tabular}{ll}
\hline
\textbf{Source Type} & \textbf{Modality} \\
\hline
Textbook & Text + Images \\
Case Report & Text + Images \\
Fact-sheet & Text \\
QA-set & Text + Images \\
\hline
\end{tabular}
\caption{Overview of Data Sources Type} 
\label{tab:data-sources-type}
\end{table}


\subsubsection{Data Splits Strategy}
% Explain how you split the data for different purposes

The collected data was strategically partitioned to support three distinct use cases:

\begin{enumerate}
    \item \textbf{RAG Knowledge Base:} 
    \begin{itemize}
        \item Purpose: External knowledge retrieval for augmented generation
        \item Composition: Medical textbooks, clinical guidelines, disease references, diagnostic criteria
        \item Size: [Specify size/percentage]
        \item Format: Chunked text with metadata (disease category, source, etc.)
    \end{itemize}
    
    \item \textbf{Fine-tuning Dataset:}
    \begin{itemize}
        \item Purpose: Model adaptation and domain-specific training
        \item Composition: Labeled clinical cases with multimodal inputs and diagnostic outputs
        \item Size: [Specify size/percentage]
        \item Format: Input-output pairs for supervised learning
    \end{itemize}
    
    \item \textbf{Test Dataset:}
    \begin{itemize}
        \item Purpose: Independent evaluation and performance assessment
        \item Composition: Held-out clinical cases unseen during training
        \item Size: [Specify size/percentage]
        \item Stratification: Balanced across disease categories and severity levels
    \end{itemize}
\end{enumerate}

% Add specific percentages based on your actual split
The data partition ratios were determined based on the following considerations:
\begin{itemize}
    \item RAG Knowledge Base: Does not require train/test separation; includes comprehensive reference material
    \item Fine-tuning Split: [e.g., 80\% train, 20\% validation] OR [e.g., 90\% train, 5\% validation, 5\% test]
    \item Test Set: Completely separate from training data to ensure unbiased evaluation
\end{itemize}

% Optional: Add figure showing data split visualization
% \begin{figure}[htbp]
%   \centering
%   \includegraphics[width=0.8\linewidth]{figures/data-splits}
%   \caption{Data split strategy for RAG, fine-tuning, and testing}
%   \label{fig:data-splits}
% \end{figure}
